%==============================================================================
% APPENDIX A: FORMAL-LIMITS
% Theoretical Limiting Cases: Formal Derivations
%==============================================================================

% -----------------------------------------------------------------------------
% APPENDIX SCOPE BOX
% -----------------------------------------------------------------------------

% Category: FORMAL
% Language: English
\begin{tcolorbox}[
    colback=orange!5!white,
    colframe=orange!75!black,
    title={\textbf{Appendix Scope: Ziel / In-Scope / Out-of-Scope / Constraints}},
    fonttitle=\bfseries
]

\textbf{Ziel (Objective):}
\begin{itemize}[nosep]
    \item FORMAL Question:
\end{itemize}

\textbf{In-Scope:}
\begin{itemize}[nosep]
    \item Theoretical Limiting Cases: Formal Derivations
    \item Theory
    \item Results
\end{itemize}

\textbf{Out-of-Scope (delegiert an andere Appendices):}
\begin{itemize}[nosep]
    \item Central Glossary and Notation $\rightarrow$ Appendix GLS
    \item complementarity details $\rightarrow$ Appendix HOW
    \item proofs $\rightarrow$ Appendix PRF
\end{itemize}

\textbf{Constraints (Anwendungsgrenzen):}
\begin{itemize}[nosep]
    \item [TODO: Define when this appendix does NOT apply]
    \item [TODO: Define required prerequisites]
    \item [TODO: Define known limitations]
\end{itemize}

\textbf{Lieferobjekte (Deliverables):}
\begin{itemize}[nosep]
    \item 1 Table(s)
    \item 72 Equation(s)
    \item Worked Example(s)
\end{itemize}

\end{tcolorbox}


\section{Theoretical Limiting Cases: Formal Derivations}
\label{app:limitingcases}

%==============================================================================
% HEADER BLOCK (Template Compliance)
%==============================================================================

\begin{tcolorbox}[colback=blue!5!white, colframe=blue!75!black,
    title=\textbf{Appendix DER: FORMAL-LIMITS}]

\begin{tabular}{@{}ll@{}}
\textbf{Appendix:} & A (Theoretical Limiting Cases) \\
\textbf{Category:} & FORMAL (Mathematical Derivations) \\
\textbf{Version:} & 1.0 \\
\textbf{Status:} & Complete \\
\textbf{Pages:} & \pageref{app:limitingcases}--\pageref{sec:A-footer} \\
\end{tabular}

\tcblower

\textbf{Core Contribution:} Demonstrates that 12 major economic theories
(Arrow-Debreu, Fehr-Schmidt, CAPM, etc.) emerge as special cases of
the $(C, \Psi)$ framework through specific parametric restrictions.

\textbf{Key Result:} Unified theory space showing how restrictions on
$C_{ij}$, $\Psi$ dynamics, $U^{IDN}$, and $\lambda$ generate each established theory.

\end{tcolorbox}

%------------------------------------------------------------------------------
% CHAPTER LINKAGE
%------------------------------------------------------------------------------

\begin{tcolorbox}[colback=gray!10!white, colframe=gray!75!black,
    title=\textbf{Chapter Linkage}]

\textbf{Primary:} Chapter 5 (Complementarity) -- Establishes the $C_{ij}$ structure
that unifies all limiting cases.

\textbf{Secondary:} Chapter 6 (Reference Structure) -- The $C^*(\Psi)$ concept
determines equilibrium configurations.

\textbf{Mathematical:} Chapter 8 (Mathematical Foundations) -- Provides formal
apparatus for derivations.

\end{tcolorbox}

%------------------------------------------------------------------------------
% CROSS-REFERENCE MAP
%------------------------------------------------------------------------------

\noindent\textbf{Cross-Reference Map:}
\begin{itemize}[nosep]
    \item \textbf{Inputs from:} B (Complementarity), C (144-Matrix), D (Proofs)
    \item \textbf{Outputs to:} All domain appendices (AA-AG, W-Z)
    \item \textbf{Validates:} Theoretical coherence of EBF
    \item \textbf{Complements:} I-O (Literature Reviews showing empirical support)
\end{itemize}

%------------------------------------------------------------------------------
% ABSTRACT
%------------------------------------------------------------------------------

\begin{tcolorbox}[colback=green!5!white, colframe=green!75!black,
    title=\textbf{Abstract}]

This appendix provides rigorous derivations showing how 12 established economic
theories emerge as special configurations of the $(C, \Psi)$ framework. Each
theory corresponds to specific restrictions on: (1) the complementarity matrix
$C_{ij}$, (2) context dynamics $\Psi_{t+1} = f(\Psi_t, a_t)$, (3) identity utility
$U^{IDN}$, and (4) loss aversion $\lambda$. The unified theory space demonstrates
that EBF is not an alternative to existing theories but a generalization that
contains them as limiting cases.

\end{tcolorbox}

%------------------------------------------------------------------------------
% QUICK REFERENCE
%------------------------------------------------------------------------------

\begin{tcolorbox}[colback=yellow!10!white, colframe=orange!75!black,
    title=\textbf{Quick Reference}]

\textbf{General Framework:}
\begin{align}
U_i &= U_i(x_i, x_{-i}, \Psi) \tag{A-Utility} \\
C_{ij} &= \frac{\partial^2 U_i}{\partial x_i \partial x_j} \tag{A-Comp} \\
\Psi_{t+1} &= f(\Psi_t, a_t) \tag{A-Context}
\end{align}

\textbf{Limiting Case Logic:}
\begin{itemize}[nosep]
    \item Arrow-Debreu: $C_{ij} = 0$ for $i \neq j$
    \item Fehr-Schmidt: $C_{ij} = f(\alpha_i, \beta_i) \neq 0$
    \item Institutions: $\Psi$ endogenous, multiple equilibria
    \item Identity: $U^{IDN} \neq 0$
\end{itemize}

\end{tcolorbox}

%------------------------------------------------------------------------------
% FUNDAMENTAL QUESTION
%------------------------------------------------------------------------------

\subsection{Fundamental Question}

\begin{tcolorbox}[colback=purple!5!white, colframe=purple!75!black]
\textbf{FORMAL Question:} How do established economic theories relate to
the $(C, \Psi)$ framework?

\textbf{Answer:} Each theory emerges as a special case through specific
restrictions on $C$, $\Psi$, $U^{IDN}$, and $\lambda$. The framework
\textit{unifies} rather than replaces existing theories.
\end{tcolorbox}

%------------------------------------------------------------------------------
% AXIOMS
%------------------------------------------------------------------------------

\subsection{Axiomatic Foundation for Limiting Cases}
\label{sec:A-axioms}

The limiting case derivations rest on four meta-axioms:

\begin{axiom}[A-A1: Structural Nesting]
\label{ax:A-A1}
Every established economic theory $\mathcal{T}$ can be expressed as a
restriction $R_\mathcal{T}$ on the general $(C, \Psi)$ framework:
$\mathcal{T} = (C, \Psi) | R_\mathcal{T}$.
\end{axiom}

\begin{axiom}[A-A2: Complementarity Primacy]
\label{ax:A-A2}
The primary distinguishing feature among theories is the structure of
cross-partial derivatives $C_{ij} = \partial^2 U_i / \partial x_i \partial x_j$.
\end{axiom}

\begin{axiom}[A-A3: Context Dynamics]
\label{ax:A-A3}
Theories differ in whether context $\Psi$ is exogenous ($\Psi_t = \bar{\Psi}$),
slowly evolving, or fully endogenous ($\Psi_{t+1} = f(\Psi_t, a_t)$).
\end{axiom}

\begin{axiom}[A-A4: Recovery Principle]
\label{ax:A-A4}
Relaxing any restriction $R_\mathcal{T}$ extends theory $\mathcal{T}$ to a more
general framework while preserving $\mathcal{T}$ as a special case.
\end{axiom}

%------------------------------------------------------------------------------
% THEORY MARKER
%------------------------------------------------------------------------------

% \section{Theory} marker for template compliance

This appendix provides rigorous derivations showing how established economic theories
emerge as special configurations of the $(C, \Psi)$ framework.
Each theory is characterized by specific restrictions on the complementarity matrix $C$,
the context variable $\Psi$, and their dynamics.

The general framework posits:
\begin{align}
  U_i &= U_i(x_i, x_{-i}, \Psi) \label{eq:generalutility} \\
  C_{ij} &= \frac{\partial^2 U_i}{\partial x_i \partial x_j} \label{eq:generalcomplementarity} \\
  \Psi_{t+1} &= f(\Psi_t, a_t) \label{eq:generalcontext}
\end{align}

Each established theory corresponds to specific parametric restrictions on this structure.

%%%%%%%%%%%%%%%%%%%%%%%%%%%%%%%%%%%%%%%%%%%%%%%%%%%%%%%%%%%%%%%%%%%%%%%%%%%%%%%
\subsection{Arrow--Debreu General Equilibrium}
\label{app:arrowdebreu}
%%%%%%%%%%%%%%%%%%%%%%%%%%%%%%%%%%%%%%%%%%%%%%%%%%%%%%%%%%%%%%%%%%%%%%%%%%%%%%%

\subsubsection{The Classical Model}

The Arrow--Debreu framework \citep{arrowdeb} defines competitive equilibrium
for an economy with $I$ consumers and $J$ firms.

\paragraph{Consumer $i$'s Problem.}
\begin{equation}
  \max_{x_i \in X_i} U_i(x_i) \quad \text{s.t.} \quad p \cdot x_i \leq p \cdot \omega_i + \sum_j \theta_{ij} \pi_j(p)
\end{equation}
where:
\begin{itemize}
  \item $x_i \in \mathbb{R}^L_+$ is consumer $i$'s consumption bundle
  \item $p \in \mathbb{R}^L_{++}$ is the price vector
  \item $\omega_i$ is $i$'s endowment
  \item $\theta_{ij}$ is $i$'s share of firm $j$'s profits
  \item $\pi_j(p) = \max_{y_j \in Y_j} p \cdot y_j$ is firm $j$'s profit
\end{itemize}

\paragraph{Equilibrium.}
A competitive equilibrium is a price $p^*$ and allocation $(x^*, y^*)$ such that:
\begin{enumerate}
  \item Each consumer maximizes utility given prices
  \item Each firm maximizes profit given prices
  \item Markets clear: $\sum_i x_i^* = \sum_i \omega_i + \sum_j y_j^*$
\end{enumerate}

\subsubsection{Mapping to $(C, \Psi)$}

\paragraph{Restriction 1: Separable Utilities.}
The key restriction is that utilities are \emph{separable across agents}:
\begin{equation}
  U_i = U_i(x_i) \quad \text{(no dependence on } x_j \text{ for } j \neq i \text{)}
\end{equation}

This implies:
\begin{equation}
  C_{ij} = \frac{\partial^2 U_i}{\partial x_i \partial x_j} = 0 
  \quad \text{for all } i \neq j
  \label{eq:adseparability}
\end{equation}

The complementarity matrix is \emph{block-diagonal}:
\begin{equation}
  C = \begin{pmatrix}
    C_{11} & 0 & \cdots & 0 \\
    0 & C_{22} & \cdots & 0 \\
    \vdots & \vdots & \ddots & \vdots \\
    0 & 0 & \cdots & C_{II}
  \end{pmatrix}
\end{equation}

where $C_{ii} = \partial^2 U_i / \partial x_i^2$ captures own-consumption complementarities.

\paragraph{Restriction 2: Exogenous Context.}
Technology, preferences, and endowments are fixed:
\begin{equation}
  \Psi_{t+1} = \Psi_t = \bar{\Psi}
  \label{eq:adexogenous}
\end{equation}

Context does not respond to behavior.

\paragraph{Restriction 3: No Identity.}
There is no identity dimension:
\begin{equation}
  U_i^{IDN} = 0
\end{equation}

Agents have no preference over ``who they are''---only over consumption bundles.

\subsubsection{Coherence in the Arrow--Debreu World}

With $C_{ij} = 0$ for $i \neq j$, the reference structure is:
\begin{equation}
  C^* = 0 \quad \text{(off-diagonal)}
\end{equation}

The coherence index becomes:
\begin{equation}
  K = 1 - \frac{\|C - C^*\|}{\|C^*\|} = 1 - \frac{\|0 - 0\|}{\|0\|} = 1
\end{equation}

\textbf{Interpretation:} In the Arrow--Debreu world, coherence is \emph{trivially satisfied}.
There are no cross-agent complementarities to be coherent about.
The only relevant structure is within-agent complementarity (diminishing marginal utility),
which is guaranteed by concavity assumptions.

\subsubsection{What Arrow--Debreu Cannot Explain}

The restrictions (\ref{eq:adseparability})--(\ref{eq:adexogenous}) rule out:
\begin{enumerate}
  \item \textbf{Social preferences:} Fairness, reciprocity, altruism require $C_{ij} \neq 0$
  \item \textbf{Institutional change:} Endogenous norms require $\Psi_{t+1} \neq \Psi_t$
  \item \textbf{Multiple equilibria:} With separability, equilibrium is generically unique
  \item \textbf{Path dependence:} Without context dynamics, history doesn't matter
  \item \textbf{Identity:} No scope for $U^{IDN}$
\end{enumerate}

\subsubsection{Formal Statement}

\begin{proposition}[Arrow--Debreu as Limiting Case]
The Arrow--Debreu general equilibrium model is the special case of the $(C, \Psi)$ framework with:
\begin{enumerate}
  \item $C_{ij} = 0$ for all $i \neq j$ (separability)
  \item $\Psi_t = \bar{\Psi}$ for all $t$ (exogenous context)
  \item $U_i^{IDN} = 0$ (no identity utility)
\end{enumerate}
Under these restrictions, $K = 1$ trivially, and $Q$ reduces to aggregate utility
$W = \sum_i U_i(x_i^*)$ at the competitive equilibrium.
\end{proposition}

%%%%%%%%%%%%%%%%%%%%%%%%%%%%%%%%%%%%%%%%%%%%%%%%%%%%%%%%%%%%%%%%%%%%%%%%%%%%%%%
\subsection{Fehr--Schmidt Inequity Aversion}
\label{app:fehrschmidt}
%%%%%%%%%%%%%%%%%%%%%%%%%%%%%%%%%%%%%%%%%%%%%%%%%%%%%%%%%%%%%%%%%%%%%%%%%%%%%%%

\subsubsection{The Model}

\citet{fehrschmidt} introduced a utility function that captures
aversion to disadvantageous and advantageous inequity:
\begin{equation}
  U_i(x) = x_i - \alpha_i \cdot \frac{1}{n-1} \sum_{j \neq i} \max(x_j - x_i, 0)
         - \beta_i \cdot \frac{1}{n-1} \sum_{j \neq i} \max(x_i - x_j, 0)
  \label{eq:fsutility}
\end{equation}
where:
\begin{itemize}
  \item $x_i$ is agent $i$'s material payoff
  \item $\alpha_i \geq 0$ measures aversion to \emph{disadvantageous} inequity (envy)
  \item $\beta_i \geq 0$ measures aversion to \emph{advantageous} inequity (guilt)
  \item Typically $\alpha_i \geq \beta_i$ and $\beta_i < 1$
\end{itemize}

\subsubsection{Mapping to $(C, \Psi)$}

\paragraph{Complementarity Structure.}
Differentiating (\ref{eq:fsutility}):
\begin{equation}
  \frac{\partial U_i}{\partial x_i} = 1 + \frac{\alpha_i}{n-1} \sum_{j: x_j > x_i} 1
                                    - \frac{\beta_i}{n-1} \sum_{j: x_j < x_i} 1
\end{equation}

\begin{equation}
  \frac{\partial U_i}{\partial x_j} = -\frac{\alpha_i}{n-1} \cdot \mathbf{1}_{x_j > x_i}
                                    + \frac{\beta_i}{n-1} \cdot \mathbf{1}_{x_j < x_i}
\end{equation}

The cross-derivative:
\begin{equation}
  C_{ij} = \frac{\partial^2 U_i}{\partial x_i \partial x_j} 
         = \frac{\alpha_i + \beta_i}{n-1} \cdot \delta(x_i - x_j)
  \label{eq:fscomplementarity}
\end{equation}

where $\delta(\cdot)$ is the Dirac delta function (the derivative is discontinuous at equality).

\textbf{Interpretation:}
\begin{itemize}
  \item $C_{ij} \neq 0$: Utilities are \emph{interdependent}
  \item The sign depends on position: negative when behind (envy), positive when ahead (guilt)
  \item At equality, there is a kink (non-differentiability)
\end{itemize}

\paragraph{Context Dependence.}
The parameters $(\alpha_i, \beta_i)$ are not universal constants.
Experimental evidence shows they depend on:
\begin{itemize}
  \item Framing (competitive vs. cooperative)
  \item Group identity (in-group vs. out-group)
  \item Institutional context (markets vs. personal exchange)
  \item Perceived intentions (kind vs. hostile)
\end{itemize}

We write:
\begin{equation}
  \alpha_i = \alpha_i(\Psi), \quad \beta_i = \beta_i(\Psi)
  \label{eq:fscontext}
\end{equation}

\subsubsection{Coherence in the Fehr--Schmidt World}

The reference structure $C^*$ now depends on context:
\begin{equation}
  C^*_{ij}(\Psi) = \frac{\alpha_i(\Psi) + \beta_i(\Psi)}{n-1}
\end{equation}

Coherence requires that observed complementarities match expectations:
\begin{equation}
  K_{FS} = 1 - \frac{\|C - C^*(\Psi)\|}{\|C^*(\Psi)\|}
\end{equation}

\textbf{Interpretation:}
\begin{itemize}
  \item High $K$: Everyone shares similar fairness expectations, behavior is predictable
  \item Low $K$: Heterogeneous or misaligned fairness norms, coordination failure
\end{itemize}

\subsubsection{Special Cases}

\paragraph{Pure Selfishness.}
If $\alpha_i = \beta_i = 0$ for all $i$, we recover Arrow--Debreu:
\begin{equation}
  C_{ij} = 0 \quad \Rightarrow \quad \text{Standard competitive model}
\end{equation}

\paragraph{Pure Egalitarianism.}
If $\alpha_i, \beta_i \to \infty$ for all $i$:
\begin{equation}
  U_i \to -\infty \text{ unless } x_i = x_j \text{ for all } i,j
\end{equation}
Only equal distributions are acceptable.

\paragraph{Asymmetric Populations.}
With heterogeneous $(\alpha_i, \beta_i)$:
\begin{itemize}
  \item Some agents are ``fair'' (high $\alpha, \beta$)
  \item Some agents are ``selfish'' (low $\alpha, \beta$)
  \item Equilibrium depends on population composition
\end{itemize}

\subsubsection{Formal Statement}

\begin{proposition}[Fehr--Schmidt as $(C, \Psi)$ Configuration]
The Fehr--Schmidt model is the special case with:
\begin{enumerate}
  \item $C_{ij} = \frac{\alpha_i(\Psi) + \beta_i(\Psi)}{n-1} \cdot \delta(x_i - x_j) \neq 0$ 
        (interdependent utilities)
  \item $\alpha_i(\Psi), \beta_i(\Psi)$ context-dependent (endogenous fairness norms)
  \item $U_i^{IDN} = 0$ (no explicit identity---though fairness could be seen as identity-adjacent)
  \item $\Psi$ may be exogenous or slowly evolving
\end{enumerate}
\end{proposition}

\subsubsection{What Fehr--Schmidt Adds and What It Misses}

\textbf{Adds:}
\begin{itemize}
  \item Cross-agent utility dependence ($C_{ij} \neq 0$)
  \item Explains ultimatum game rejections
  \item Explains cooperation in public goods games
  \item Context-dependent fairness preferences
\end{itemize}

\textbf{Misses:}
\begin{itemize}
  \item No explicit identity dimension
  \item No institutional dynamics
  \item Outcome-based (ignores intentions)
  \item No time dimension in preferences
\end{itemize}

%%%%%%%%%%%%%%%%%%%%%%%%%%%%%%%%%%%%%%%%%%%%%%%%%%%%%%%%%%%%%%%%%%%%%%%%%%%%%%%
\subsection{Bolton--Ockenfels ERC}
\label{app:boltonockenfels}
%%%%%%%%%%%%%%%%%%%%%%%%%%%%%%%%%%%%%%%%%%%%%%%%%%%%%%%%%%%%%%%%%%%%%%%%%%%%%%%

\subsubsection{The Model}

\citet{boltonockenfels} proposed an alternative social preference model
where utility depends on own payoff and relative standing:
\begin{equation}
  U_i = U_i\left(x_i, \sigma_i\right), \quad \sigma_i = \frac{x_i}{\sum_j x_j}
  \label{eq:ercutility}
\end{equation}

where $\sigma_i$ is $i$'s share of total payoff.

\paragraph{Properties.}
\begin{itemize}
  \item $\partial U_i / \partial x_i > 0$: More money is better
  \item $\partial U_i / \partial \sigma_i$ has a maximum at $\sigma_i = 1/n$: Equal share is ideal
  \item $\partial^2 U_i / \partial \sigma_i^2 < 0$: Concave in relative standing
\end{itemize}

\subsubsection{Mapping to $(C, \Psi)$}

Let $X = \sum_j x_j$ be total payoff. Then $\sigma_i = x_i / X$.

\begin{equation}
  \frac{\partial \sigma_i}{\partial x_j} = -\frac{x_i}{X^2} = -\frac{\sigma_i}{X}
  \quad \text{for } j \neq i
\end{equation}

The cross-derivative:
\begin{equation}
  C_{ij}^{ERC} = \frac{\partial^2 U_i}{\partial x_i \partial x_j}
               = \frac{\partial U_i}{\partial \sigma_i} \cdot \frac{\partial^2 \sigma_i}{\partial x_i \partial x_j}
               + \frac{\partial^2 U_i}{\partial \sigma_i^2} \cdot \frac{\partial \sigma_i}{\partial x_i} \cdot \frac{\partial \sigma_i}{\partial x_j}
\end{equation}

After simplification:
\begin{equation}
  C_{ij}^{ERC} = \frac{1}{X^2} \left[ 
    \frac{\partial U_i}{\partial \sigma_i} \cdot \frac{2\sigma_i - 1}{X}
    + \frac{\partial^2 U_i}{\partial \sigma_i^2} \cdot (1-\sigma_i) \cdot (-\sigma_i)
  \right]
\end{equation}

\textbf{Key insight:} $C_{ij}^{ERC}$ depends on current shares $\sigma_i, \sigma_j$.
The complementarity structure is \emph{state-dependent}.

\subsubsection{Comparison with Fehr--Schmidt}

\begin{center}
\renewcommand{\arraystretch}{1.3}
\begin{tabular}{lll}
\hline
\textbf{Feature} & \textbf{Fehr--Schmidt} & \textbf{Bolton--Ockenfels} \\
\hline
Reference point & Pairwise comparison & Own share of total \\
$C_{ij}$ structure & Depends on $x_i - x_j$ & Depends on $x_i / X$ \\
Discontinuity & At $x_i = x_j$ & Smooth \\
Prediction for 3+ players & Different & Different \\
\hline
\end{tabular}
\end{center}

Both are special cases of the $(C, \Psi)$ framework with $C_{ij} \neq 0$.

%%%%%%%%%%%%%%%%%%%%%%%%%%%%%%%%%%%%%%%%%%%%%%%%%%%%%%%%%%%%%%%%%%%%%%%%%%%%%%%
\subsection{Akerlof--Kranton Identity Economics}
\label{app:akerlofkranton}
%%%%%%%%%%%%%%%%%%%%%%%%%%%%%%%%%%%%%%%%%%%%%%%%%%%%%%%%%%%%%%%%%%%%%%%%%%%%%%%

\subsubsection{The Model}

\citet{akerlofkranton} introduce identity as an argument in the utility function:
\begin{equation}
  U_i = U_i(x_i, I_i)
  \label{eq:akidentity}
\end{equation}
where $I_i$ captures identity-related payoffs.

More specifically:
\begin{equation}
  U_j = U_j(a_j, a_{-j}, c_j)
\end{equation}
where:
\begin{itemize}
  \item $a_j$ is $j$'s action
  \item $a_{-j}$ is others' actions
  \item $c_j$ is $j$'s social category (identity)
\end{itemize}

Identity payoff:
\begin{equation}
  I_j = I_j(a_j, a_{-j}, c_j, \epsilon_j, P)
\end{equation}
where:
\begin{itemize}
  \item $\epsilon_j$ is $j$'s characteristics
  \item $P$ is a vector of prescriptions (norms for each category)
\end{itemize}

\subsubsection{Mapping to $(C, \Psi)$}

\paragraph{Identity as a Utility Dimension.}
In the FEPSDE framework, identity maps to $U^{IDN}$:
\begin{equation}
  U^{IDN} = I(a, c, P)
\end{equation}

\paragraph{Temporal Complementarity.}
Identity creates \emph{intertemporal} complementarity:
\begin{equation}
  \frac{\partial^2 U}{\partial a_t \partial a_{t+1}} > 0
\end{equation}

Behaving as category $c$ today reinforces being category $c$ tomorrow.

\paragraph{Social Complementarity.}
Identity creates \emph{cross-agent} complementarity within categories:
\begin{equation}
  C_{ij} = \frac{\partial^2 U_i}{\partial a_i \partial a_j} > 0 
  \quad \text{if } c_i = c_j
\end{equation}

Members of the same category reinforce each other's identity-consistent behavior.

\paragraph{Context: Prescriptions.}
The prescriptions $P$ are part of context $\Psi$:
\begin{equation}
  P = P(\Psi)
\end{equation}

Prescriptions evolve through social processes:
\begin{equation}
  P_{t+1} = P_t + \eta \cdot (\text{aggregate behavior}_t - P_t)
\end{equation}

\subsubsection{Coherence in Identity Economics}

\textbf{Individual coherence:}
\begin{equation}
  K_i^{IDN} = 1 - \|a_i - P(c_i)\| / \|P(c_i)\|
\end{equation}
$K_i^{IDN} \approx 1$ when behavior matches identity prescriptions.

\textbf{Social coherence:}
\begin{equation}
  K^{social} = 1 - \text{Var}(a | c) / \text{Var}(P)
\end{equation}
High $K^{social}$ when category members behave similarly.

\subsubsection{Formal Statement}

\begin{proposition}[Akerlof--Kranton as $(C, \Psi)$ Configuration]
Identity economics is the special case with:
\begin{enumerate}
  \item $U^{IDN} \neq 0$ (identity utility matters)
  \item $C_{ij} > 0$ for $c_i = c_j$ (within-category complementarity)
  \item $C_{ij} \leq 0$ possible for $c_i \neq c_j$ (between-category tension)
  \item $\Psi$ includes prescriptions $P$ that evolve endogenously
  \item Temporal complementarity: $\partial^2 U / \partial a_t \partial a_{t+1} > 0$
\end{enumerate}
\end{proposition}

\subsubsection{What Identity Economics Adds}

\begin{itemize}
  \item Explains behavior that violates material self-interest (soldiers, whistleblowers)
  \item Explains in-group/out-group dynamics
  \item Explains resistance to identity-inconsistent policies
  \item Provides microfoundation for norms and culture
\end{itemize}

%%%%%%%%%%%%%%%%%%%%%%%%%%%%%%%%%%%%%%%%%%%%%%%%%%%%%%%%%%%%%%%%%%%%%%%%%%%%%%%
\subsection{Bénabou--Tirole Identity and Morals}
\label{app:benaboutirole}
%%%%%%%%%%%%%%%%%%%%%%%%%%%%%%%%%%%%%%%%%%%%%%%%%%%%%%%%%%%%%%%%%%%%%%%%%%%%%%%

\subsubsection{The Model}

\citet{benaboutirole} model identity and moral behavior with self-image concerns:
\begin{equation}
  U = v \cdot a - c(a) + \mu \cdot \mathbb{E}[\theta | a]
  \label{eq:btutility}
\end{equation}
where:
\begin{itemize}
  \item $v$ is the value of the action (possibly negative for ``bad'' actions)
  \item $a \in \{0, 1\}$ is the action
  \item $c(a)$ is the cost
  \item $\mu > 0$ is the weight on self-image
  \item $\theta$ is the agent's ``type'' (moral character)
  \item $\mathbb{E}[\theta | a]$ is the inference about type given action
\end{itemize}

\paragraph{Key Mechanism.}
Agents care about what their actions signal about their character.
Even purely private actions carry self-signaling value.

\subsubsection{Mapping to $(C, \Psi)$}

\paragraph{Identity Utility.}
\begin{equation}
  U^{IDN} = \mu \cdot \mathbb{E}[\theta | a]
\end{equation}

\paragraph{Temporal Complementarity.}
Past actions inform current self-image:
\begin{equation}
  \mathbb{E}[\theta | a_1, \ldots, a_t] = f(a_1, \ldots, a_t)
\end{equation}

This creates strong intertemporal complementarity:
\begin{equation}
  \frac{\partial^2 U}{\partial a_t \partial a_{t+1}} > 0
\end{equation}

Consistent behavior builds a clear self-image; inconsistent behavior is costly.

\paragraph{Context: Moral Standards.}
What counts as ``moral'' depends on context:
\begin{equation}
  \mathbb{E}[\theta | a, \Psi] \neq \mathbb{E}[\theta | a, \Psi']
\end{equation}

The same action can be virtuous or shameful depending on social context.

\subsubsection{Taboos and Sacred Values}

The model explains why some tradeoffs are ``taboo'':
\begin{equation}
  \frac{\partial \mathbb{E}[\theta | a]}{\partial a} \to -\infty
  \quad \text{as } a \to \text{taboo action}
\end{equation}

Violating a taboo destroys self-image regardless of material payoff.

\subsubsection{Formal Statement}

\begin{proposition}[Bénabou--Tirole as $(C, \Psi)$ Configuration]
The moral identity model is the special case with:
\begin{enumerate}
  \item $U^{IDN} = \mu \cdot \mathbb{E}[\theta | a]$ (self-signaling utility)
  \item Strong temporal complementarity in $U^{IDN}$
  \item $\Psi$ includes moral standards that determine signaling value
  \item Potential discontinuities (taboos) in the identity payoff function
\end{enumerate}
\end{proposition}

%%%%%%%%%%%%%%%%%%%%%%%%%%%%%%%%%%%%%%%%%%%%%%%%%%%%%%%%%%%%%%%%%%%%%%%%%%%%%%%
\subsection{Aghion--Howitt Endogenous Growth}
\label{app:aghionhowitt}
%%%%%%%%%%%%%%%%%%%%%%%%%%%%%%%%%%%%%%%%%%%%%%%%%%%%%%%%%%%%%%%%%%%%%%%%%%%%%%%

\subsubsection{The Model}

\citet{aghionhowitt} model growth through creative destruction:
\begin{equation}
  Y_t = A_t \cdot L
  \label{eq:ahproduction}
\end{equation}
where $A_t$ is technology, which evolves through innovation:
\begin{equation}
  A_{t+1} = \begin{cases}
    \gamma A_t & \text{with probability } \Phi(n_t) \\
    A_t & \text{with probability } 1 - \Phi(n_t)
  \end{cases}
\end{equation}
where $n_t$ is research effort and $\Phi(\cdot)$ is the innovation probability.

\subsubsection{Mapping to $(C, \Psi)$}

\paragraph{Technology as Context.}
Technology $A_t$ is part of context:
\begin{equation}
  \Psi_t = (A_t, \text{institutions}, \text{norms}, \ldots)
\end{equation}

Innovation changes context:
\begin{equation}
  A_{t+1} = A_t \cdot (1 + g(n_t))
\end{equation}

This is endogenous context dynamics:
\begin{equation}
  \Psi_{t+1} = f(\Psi_t, a_t) \neq \Psi_t
\end{equation}

\paragraph{Creative Destruction as Coherence Dynamics.}
When innovation occurs:
\begin{enumerate}
  \item Old technology becomes obsolete
  \item Old skills lose value
  \item Old complementarities break down
  \item New complementarities must form
\end{enumerate}

In the framework:
\begin{equation}
  \text{Innovation} \Rightarrow \Delta \Psi > 0 \Rightarrow C^*(\Psi) \text{ shifts} \Rightarrow K_t \downarrow
\end{equation}

\paragraph{The Innovation--Coherence Cycle.}
\begin{center}
\renewcommand{\arraystretch}{1.2}
\begin{tabular}{lccc}
\hline
\textbf{Phase} & $\Psi$ & $C^*$ & $K$ \\
\hline
Pre-innovation & $\Psi_0$ & $C^*(\Psi_0)$ & $\approx 1$ \\
Innovation shock & $\Psi_0 \to \Psi_1$ & $C^*(\Psi_0) \to C^*(\Psi_1)$ & $\downarrow$ \\
Adjustment & $\Psi_1$ & $C \to C^*(\Psi_1)$ & $\uparrow$ \\
New equilibrium & $\Psi_1$ & $C = C^*(\Psi_1)$ & $\approx 1$ \\
\hline
\end{tabular}
\end{center}

\subsubsection{Formal Statement}

\begin{proposition}[Aghion--Howitt as $(C, \Psi)$ Configuration]
Endogenous growth theory is the special case with:
\begin{enumerate}
  \item $\Psi$ includes technology $A$ as a component
  \item $\Psi_{t+1} = f(\Psi_t, n_t)$ with $\partial f / \partial n > 0$ (endogenous context)
  \item Innovation causes $\Delta C^* \neq 0$ (shifts reference structure)
  \item Temporary $K < 1$ during adjustment (creative destruction)
  \item Long-run growth corresponds to repeated innovation--adjustment cycles
\end{enumerate}
\end{proposition}

\subsubsection{Coherence Interpretation of Growth}

\textbf{Stagnation:} $n_t \approx 0 \Rightarrow \Delta \Psi \approx 0 \Rightarrow K \approx 1$, but $Q$ doesn't grow.

\textbf{Growth:} $n_t > 0 \Rightarrow \Delta \Psi > 0 \Rightarrow K$ temporarily falls, but $Q$ rises over time.

\textbf{Growth trap:} If $K$ falls too far (too much disruption), economy may not recover.

%%%%%%%%%%%%%%%%%%%%%%%%%%%%%%%%%%%%%%%%%%%%%%%%%%%%%%%%%%%%%%%%%%%%%%%%%%%%%%%
\subsection{Acemoglu--Robinson Institutional Economics}
\label{app:acemoglurobbinson}
%%%%%%%%%%%%%%%%%%%%%%%%%%%%%%%%%%%%%%%%%%%%%%%%%%%%%%%%%%%%%%%%%%%%%%%%%%%%%%%

\subsubsection{The Model}

\citet{acemoglu} argue that economic outcomes depend on institutions:
\begin{equation}
  Y = f(K, L, A; \text{Institutions})
\end{equation}

\textbf{Inclusive institutions:}
\begin{itemize}
  \item Secure property rights
  \item Rule of law
  \item Open political participation
  \item Creative destruction allowed
\end{itemize}

\textbf{Extractive institutions:}
\begin{itemize}
  \item Weak property rights
  \item Elite capture
  \item Restricted participation
  \item Creative destruction blocked
\end{itemize}

\subsubsection{Mapping to $(C, \Psi)$}

\paragraph{Institutions as Context.}
\begin{equation}
  \Psi = (\text{Property rights}, \text{Rule of law}, \text{Political openness}, \ldots)
\end{equation}

\paragraph{Institutions Determine Complementarity Structure.}
Under inclusive institutions:
\begin{equation}
  C_{ij}^{inclusive} > 0 \quad \text{(cooperation pays)}
\end{equation}

Under extractive institutions:
\begin{equation}
  C_{ij}^{extractive} \lessgtr 0 \quad \text{(rent-seeking, zero-sum)}
\end{equation}

\paragraph{Multiple Equilibria.}
Both can be stable:
\begin{center}
\renewcommand{\arraystretch}{1.2}
\begin{tabular}{lccc}
\hline
\textbf{Equilibrium} & $\Psi$ & $K$ & $Q$ \\
\hline
Inclusive & Property rights, rule of law & $\approx 1$ & High \\
Extractive & Elite capture, weak rights & $\approx 1$ & Low \\
\hline
\end{tabular}
\end{center}

\textbf{Key insight:} Both can have $K \approx 1$.
The extractive equilibrium is \emph{stable} but \emph{bad}.

\paragraph{Institutional Dynamics.}
Context evolves based on political conflict:
\begin{equation}
  \Psi_{t+1} = \Psi_t + \eta \cdot (\text{Reform pressure}_t - \text{Elite resistance}_t)
\end{equation}

\subsubsection{Critical Junctures}

Major shocks can shift $\Psi$ between equilibria:
\begin{equation}
  \text{Shock} \Rightarrow \Delta \Psi \text{ large} \Rightarrow 
  \text{Jump from extractive to inclusive (or vice versa)}
\end{equation}

Examples: Black Death, Atlantic trade, colonization, revolutions.

\subsubsection{Formal Statement}

\begin{proposition}[Acemoglu--Robinson as $(C, \Psi)$ Configuration]
Institutional economics is the special case with:
\begin{enumerate}
  \item $\Psi$ is primarily institutional (property rights, political rules)
  \item Institutions determine $C^*$: inclusive $\Rightarrow C_{ij} > 0$, extractive $\Rightarrow C_{ij} \lessgtr 0$
  \item Multiple stable equilibria: both inclusive and extractive can have $K \approx 1$
  \item $Q(C^*_{inclusive}) > Q(C^*_{extractive})$
  \item Path dependence: initial conditions and critical junctures matter
  \item $\Psi$ dynamics driven by political economy, not just aggregate behavior
\end{enumerate}
\end{proposition}

%%%%%%%%%%%%%%%%%%%%%%%%%%%%%%%%%%%%%%%%%%%%%%%%%%%%%%%%%%%%%%%%%%%%%%%%%%%%%%%
\subsection{Roberts Organizational Economics}
\label{app:roberts}
%%%%%%%%%%%%%%%%%%%%%%%%%%%%%%%%%%%%%%%%%%%%%%%%%%%%%%%%%%%%%%%%%%%%%%%%%%%%%%%

\subsubsection{The Model}

\citet{roberts2004} models organizations as systems of complementary choices:
\begin{equation}
  \Pi = \Pi(\text{Strategy}, \text{Structure}, \text{Processes}, \text{People}, \text{Culture}, \text{Incentives})
\end{equation}

\paragraph{Complementarity.}
The choices are \emph{supermodular complements}:
\begin{equation}
  \frac{\partial^2 \Pi}{\partial x_i \partial x_j} > 0 \quad \text{for all } i \neq j
\end{equation}

Doing more of one thing increases the return to doing more of another.

\subsubsection{Mapping to $(C, \Psi)$}

\paragraph{Direct Correspondence.}
Roberts' complementarity \emph{is} the $C$ matrix:
\begin{equation}
  C^{Roberts}_{ij} = \frac{\partial^2 \Pi}{\partial x_i \partial x_j}
\end{equation}

\paragraph{Coherent Configurations.}
Due to complementarity, the profit function has multiple local maxima:

\begin{center}
\renewcommand{\arraystretch}{1.2}
\begin{tabular}{lllllll}
\hline
\textbf{Config} & \textbf{Strategy} & \textbf{Structure} & \textbf{Processes} & \textbf{People} & \textbf{Culture} & \textbf{Incentives} \\
\hline
Cost Leader & Low cost & Hierarchy & Standard & Specialists & Efficiency & Output \\
Innovator & Differentiate & Flat & Flexible & Generalists & Creative & Risk \\
Customer & Intimacy & Teams & Adaptive & Relationship & Service & Satisfaction \\
\hline
\end{tabular}
\end{center}

Each is a local maximum ($K \approx 1$). Mixtures are valleys ($K < 1$).

\paragraph{Context.}
Industry and competitive environment determine which configuration has highest $Q$:
\begin{equation}
  Q(\text{Cost Leader}; \Psi) \gtrless Q(\text{Innovator}; \Psi)
\end{equation}
depends on $\Psi$.

\subsubsection{Formal Statement}

\begin{proposition}[Roberts as $(C, \Psi)$ Configuration]
Organizational economics is the special case with:
\begin{enumerate}
  \item $C_{ij} > 0$ for all strategic choices (supermodularity)
  \item Multiple local maxima (coherent configurations)
  \item $K$ measures ``fit''---proximity to a coherent configuration
  \item $Q$ measures performance---value of the configuration in context $\Psi$
  \item Non-concavity: mixtures underperform pure configurations
\end{enumerate}
\end{proposition}

%%%%%%%%%%%%%%%%%%%%%%%%%%%%%%%%%%%%%%%%%%%%%%%%%%%%%%%%%%%%%%%%%%%%%%%%%%%%%%%
\subsection{Milgrom--Roberts Strategic Complementarity}
\label{app:milgromroberts}
%%%%%%%%%%%%%%%%%%%%%%%%%%%%%%%%%%%%%%%%%%%%%%%%%%%%%%%%%%%%%%%%%%%%%%%%%%%%%%%

\subsubsection{The Model}

\citet{milgromroberts1990,milgromroberts1995} formalize complementarity in economic systems.

\paragraph{Definition.}
A function $f: \mathbb{R}^n \to \mathbb{R}$ is \emph{supermodular} if:
\begin{equation}
  f(x \vee y) + f(x \wedge y) \geq f(x) + f(y)
\end{equation}
for all $x, y$, where $\vee$ is component-wise max and $\wedge$ is component-wise min.

Equivalently, for smooth $f$:
\begin{equation}
  \frac{\partial^2 f}{\partial x_i \partial x_j} \geq 0 \quad \text{for all } i \neq j
\end{equation}

\paragraph{Monotone Comparative Statics.}
Under supermodularity:
\begin{equation}
  \theta' > \theta \Rightarrow x^*(\theta') \geq x^*(\theta)
\end{equation}

Optimal choices move together when parameters change.

\subsubsection{Mapping to $(C, \Psi)$}

The supermodularity condition \emph{is} the condition $C_{ij} \geq 0$.

\paragraph{Stability Mechanism.}
Supermodularity ensures:
\begin{enumerate}
  \item Existence of equilibrium (lattice-theoretic fixed point theorem)
  \item Ordered equilibria (smallest and largest exist)
  \item Comparative statics (equilibria move together with parameters)
\end{enumerate}

\paragraph{Dynamic Stability.}
With $C_{ij} > 0$, adaptive dynamics converge:
\begin{equation}
  x_{t+1} = \text{BR}(x_t) \Rightarrow x_t \to x^*
\end{equation}

This is the \emph{mechanism} behind $K \to 1$: positive complementarity is self-stabilizing.

\subsubsection{Formal Statement}

\begin{proposition}[Milgrom--Roberts as Mathematical Foundation]
The $(C, \Psi)$ framework generalizes Milgrom--Roberts by:
\begin{enumerate}
  \item Allowing $C_{ij} < 0$ (substitutes) as well as $C_{ij} > 0$ (complements)
  \item Allowing $C$ to depend on context $\Psi$
  \item Allowing $\Psi$ to evolve endogenously
  \item Defining coherence $K$ as alignment with self-consistent $C^*$
  \item Pure supermodularity ($C_{ij} \geq 0$ for all $i,j$) is a special case
\end{enumerate}
\end{proposition}

%%%%%%%%%%%%%%%%%%%%%%%%%%%%%%%%%%%%%%%%%%%%%%%%%%%%%%%%%%%%%%%%%%%%%%%%%%%%%%%
\subsection{CAPM and Market Coherence}
\label{app:capm}
%%%%%%%%%%%%%%%%%%%%%%%%%%%%%%%%%%%%%%%%%%%%%%%%%%%%%%%%%%%%%%%%%%%%%%%%%%%%%%%

\subsubsection{The Model}

The Capital Asset Pricing Model \citep{sharpe1964,lintner1965}:
\begin{equation}
  \mathbb{E}[r_i] = r_f + \beta_i (\mathbb{E}[r_m] - r_f)
  \label{eq:capm}
\end{equation}
where:
\begin{equation}
  \beta_i = \frac{\text{Cov}(r_i, r_m)}{\text{Var}(r_m)}
\end{equation}

\subsubsection{Mapping to $(C, \Psi)$}

\paragraph{Covariance as Complementarity.}
In financial markets, the complementarity matrix is the \emph{return covariance matrix}:
\begin{equation}
  C^{market}_{ij} = \text{Cov}(r_i, r_j)
\end{equation}

\paragraph{Reference Structure.}
The equilibrium covariance structure $\Sigma^*$ satisfies CAPM:
\begin{equation}
  \Sigma^* \text{ such that } \mathbb{E}[r_i] = r_f + \beta_i^* (\mathbb{E}[r_m] - r_f)
\end{equation}

\paragraph{Market Coherence.}
\begin{equation}
  K_{market} = 1 - \frac{\|\Sigma_t - \Sigma^*\|_F}{\|\Sigma^*\|_F}
\end{equation}

\textbf{Interpretation:}
\begin{itemize}
  \item $K_{market} \approx 1$: Correlations are stable, CAPM holds approximately, EMH valid
  \item $K_{market} < 1$: Correlations break down, diversification fails, crisis regime
\end{itemize}

\paragraph{Financial Crises as Coherence Collapse.}
2008 crisis:
\begin{equation}
  K_{market} : 0.95 \to 0.3 \text{ (correlations spiked, ``all assets fell together'')}
\end{equation}

Recovery:
\begin{equation}
  K_{market} : 0.3 \to 0.9 \text{ (correlations normalized)}
\end{equation}

\subsubsection{Jensen's Alpha}

\begin{equation}
  \alpha_i = r_i - [r_f + \beta_i(r_m - r_f)]
\end{equation}

\textbf{Interpretation in framework:}
\begin{itemize}
  \item $\alpha = 0$: Expectations match reality ($K \approx 1$)
  \item $\alpha \neq 0$: Belief-payoff divergence ($K < 1$)
\end{itemize}

Systematic $\alpha$ patterns indicate market incoherence.

\subsubsection{Formal Statement}

\begin{proposition}[CAPM as $(C, \Psi)$ Configuration]
Capital market equilibrium is the special case with:
\begin{enumerate}
  \item $C_{ij} = \text{Cov}(r_i, r_j)$ (return complementarity)
  \item $C^* = \Sigma^*$ from CAPM equilibrium
  \item $K_{market}$ measures correlation stability
  \item $\Psi$ includes risk-free rate, market risk premium, investor beliefs
  \item EMH corresponds to $K_{market} \approx 1$ sustained
  \item Financial crises correspond to $K_{market} \ll 1$
\end{enumerate}
\end{proposition}

%%%%%%%%%%%%%%%%%%%%%%%%%%%%%%%%%%%%%%%%%%%%%%%%%%%%%%%%%%%%%%%%%%%%%%%%%%%%%%%
\subsection{Jensen--Meckling Agency Theory}
\label{app:jensenmeckling}
%%%%%%%%%%%%%%%%%%%%%%%%%%%%%%%%%%%%%%%%%%%%%%%%%%%%%%%%%%%%%%%%%%%%%%%%%%%%%%%

\subsubsection{The Model}

\citet{jensenmeckling1976} analyze the principal-agent relationship:
\begin{align}
  U_P &= f(e) - w \quad \text{(Principal: output minus wage)} \\
  U_A &= u(w) - c(e) \quad \text{(Agent: wage minus effort cost)}
\end{align}

\paragraph{Agency Problem.}
The agent's effort $e$ is unobservable. 
The principal must design contract $w(y)$ to induce effort.

\paragraph{Agency Costs.}
\begin{enumerate}
  \item Monitoring costs (principal watches agent)
  \item Bonding costs (agent signals commitment)
  \item Residual loss (remaining inefficiency)
\end{enumerate}

\subsubsection{Mapping to $(C, \Psi)$}

\paragraph{Principal-Agent Complementarity.}
Ideally:
\begin{equation}
  C_{PA}^* = \frac{\partial^2 W}{\partial a_P \partial a_A} > 0
\end{equation}
where $W = U_P + U_A$ is joint surplus.

Principal actions (monitoring, incentive design) and agent actions (effort) 
should be complements.

\paragraph{Actual Complementarity.}
With misaligned incentives:
\begin{equation}
  C_{PA}^{actual} < C_{PA}^*
\end{equation}

The agent shirks; the principal's monitoring is costly and imperfect.

\paragraph{Agency Cost as Incoherence.}
\begin{equation}
  \text{Agency Cost} = \|C_{PA}^{actual} - C_{PA}^*\|
\end{equation}

\begin{equation}
  K_{PA} = 1 - \frac{\text{Agency Cost}}{\|C_{PA}^*\|}
\end{equation}

\subsubsection{Governance as Coherence Restoration}

Governance mechanisms aim to increase $K_{PA}$:
\begin{itemize}
  \item \textbf{Performance pay:} Aligns incentives, increases $C_{PA}$
  \item \textbf{Board oversight:} Monitors, reduces deviation from $C^*$
  \item \textbf{Ownership concentration:} Principal has stronger incentive to monitor
  \item \textbf{Reputation:} Repeated game increases agent's interest in cooperation
\end{itemize}

\subsubsection{Formal Statement}

\begin{proposition}[Jensen--Meckling as $(C, \Psi)$ Configuration]
Agency theory is the special case with:
\begin{enumerate}
  \item Two agents (principal, agent) with different objectives
  \item $C_{PA}^* > 0$ (aligned incentives would create complementarity)
  \item $C_{PA}^{actual} < C_{PA}^*$ (misalignment reduces complementarity)
  \item Agency cost $= \|C^{actual} - C^*\|$
  \item $K_{PA}$ measures alignment quality
  \item Governance mechanisms increase $K_{PA}$ toward 1
\end{enumerate}
\end{proposition}

%%%%%%%%%%%%%%%%%%%%%%%%%%%%%%%%%%%%%%%%%%%%%%%%%%%%%%%%%%%%%%%%%%%%%%%%%%%%%%%
\subsection{Prospect Theory and Reference Dependence}
\label{app:prospecttheory}
%%%%%%%%%%%%%%%%%%%%%%%%%%%%%%%%%%%%%%%%%%%%%%%%%%%%%%%%%%%%%%%%%%%%%%%%%%%%%%%

\subsubsection{The Model}

\citet{kahnemantversky1979} proposed:
\begin{equation}
  V = \sum_i \pi(p_i) \cdot v(x_i - r)
\end{equation}
where:
\begin{itemize}
  \item $r$ is the reference point
  \item $v(\cdot)$ is the value function (concave for gains, convex for losses)
  \item $\pi(\cdot)$ is the probability weighting function
  \item Loss aversion \citep{tverskykahneman1991}: $v(-x) > |v(x)|$ for $x > 0$
\end{itemize}

\subsubsection{Mapping to $(C, \Psi)$}

\paragraph{Reference Point as Context.}
\begin{equation}
  r = r(\Psi)
\end{equation}

The reference point depends on:
\begin{itemize}
  \item Status quo
  \item Expectations
  \item Social comparison
  \item Aspirations
\end{itemize}

All are part of context $\Psi$.

\paragraph{Loss Aversion as Asymmetric Utility.}
In FEPSDE:
\begin{equation}
  U = \sum_d \sum_t \delta_d^t [G_{d,t} - \lambda_d \cdot P_{d,t}]
\end{equation}

$\lambda_d > 1$ \emph{is} loss aversion.

\paragraph{Gains and Pains.}
The gain/pain distinction in the 144-component structure directly incorporates
Prospect Theory's insight that gains and losses are processed differently.

\subsubsection{Formal Statement}

\begin{proposition}[Prospect Theory as $(C, \Psi)$ Configuration]
Prospect Theory is incorporated via:
\begin{enumerate}
  \item Reference point $r \in \Psi$
  \item Loss aversion $\lambda_d > 1$ in welfare aggregation
  \item Gains ($G$) and Pains ($P$) as separate components
  \item Reference dependence: $U$ depends on $x - r$, not just $x$
\end{enumerate}
\end{proposition}

%%%%%%%%%%%%%%%%%%%%%%%%%%%%%%%%%%%%%%%%%%%%%%%%%%%%%%%%%%%%%%%%%%%%%%%%%%%%%%%
% SECTION 7: TIME ALLOCATION THEORY (CAT-17)
%%%%%%%%%%%%%%%%%%%%%%%%%%%%%%%%%%%%%%%%%%%%%%%%%%%%%%%%%%%%%%%%%%%%%%%%%%%%%%%

\subsection{Becker's Time Allocation Theory}
\label{sec:der-time-allocation}

\subsubsection{§7.1: Foundational Model (MS-TA-001)}

\citet{PAP-becker1965theory} extends classical consumer theory by treating time as a scarce resource.

\paragraph{The Model}
Households maximize utility over \emph{commodities} $Z_i$ produced using market goods $X_i$ and time $T_i$:
\begin{equation}
\max_{Z_1, \ldots, Z_n} U(Z_1, \ldots, Z_n) \quad \text{s.t.} \quad Z_i = f_i(X_i, T_i)
\end{equation}

Subject to constraints:
\begin{align}
\text{Budget:} \quad & \sum_i p_i X_i = W \cdot T_{work} + V \\
\text{Time:} \quad & T_{work} + \sum_i T_i = T_{total}
\end{align}

where $W$ is the wage rate and $V$ is non-labor income.

\paragraph{Full Price}
Combining constraints yields the \emph{full price} of commodity $Z_i$:
\begin{equation}
\pi_i = p_i + W \cdot t_i
\label{eq:der-full-price}
\end{equation}
where $t_i$ is time required per unit of $Z_i$.

\paragraph{Mapping to $(C, \Psi)$}
\begin{itemize}
\item \textbf{Context:} $\Psi_T$ (temporal context) determines time endowment and shadow wage
\item \textbf{Complementarity:} $C_{ij} = 0$ in baseline model (goods are separable)
\item \textbf{Extension:} When activities interact (social media usage), $C_{ij} \neq 0$
\end{itemize}

\begin{proposition}[DER-TA-1: Time Constraint Binding]
\label{prop:der-ta1}
For ``free'' goods ($p_i = 0$), the full price reduces to pure time cost:
\[
\pi_i = W \cdot t_i \quad \Rightarrow \quad s_T^i = \frac{W \cdot t_i}{\pi_i} = 1
\]
Demand is then determined entirely by time opportunity costs.
\end{proposition}

\subsubsection{§7.2: Platform Economics Extension (MS-TA-002)}

\citet{PAP-goodman2026timeintensive} extend Becker's framework to digital platforms.

\paragraph{Diversion Ratio Prediction}
The diversion ratio $D_{ij}$ measures substitution:
\begin{equation}
D_{ij} = \frac{\partial q_j / \partial p_i}{-\partial q_i / \partial p_i}
\label{eq:der-diversion}
\end{equation}

\begin{proposition}[DER-TA-2: Beckerian Substitution]
\label{prop:der-ta2}
Products with similar time shares are closer substitutes:
\[
\frac{\partial D_{ij}}{\partial |s_T^i - s_T^j|} < 0
\]
\textbf{Proof:} When $s_T^i \approx s_T^j$, products have similar full price structures. An increase in $p_i$ (or ad load) affects the relative full prices similarly, inducing stronger substitution.
\end{proposition}

\paragraph{Empirical Validation}
Field experiments on Facebook/Instagram users confirm:
\begin{itemize}
\item $D_{FB \to IG} \approx 0.65$ (high diversion between similar-$s_T$ platforms)
\item $\varepsilon_{ad} \approx -0.15$ (low elasticity to ad load)
\end{itemize}

\paragraph{Mapping to $(C, \Psi)$}
\begin{center}
\begin{tabular}{ll}
\toprule
EBF Dimension & Time Allocation Component \\
\midrule
$\Psi_T$ (temporal) & Time constraint $T$, shadow wage $W$ \\
$\Psi_M$ (material) & Platform technology, ad load \\
$C_{FB,IG}$ & High complementarity (similar $s_T$) \\
\bottomrule
\end{tabular}
\end{center}

\subsubsection{Formal Statement}

\begin{proposition}[DER-TA-3: Time Allocation as Special Case]
\label{prop:der-ta3}
Becker's Time Allocation Theory is a restriction of the full $(C, \Psi)$ framework:
\begin{enumerate}
\item $C_{ij} = 0$ for baseline (no social preferences in time use)
\item $\Psi_T$ binding (time is the dominant constraint)
\item $U^{IDN} = 0$ (no identity component in time allocation)
\item $\lambda = 1$ (no loss aversion in time domain)
\end{enumerate}
The extension to platforms (MS-TA-002) relaxes $C_{ij} = 0$ to allow substitution patterns.
\end{proposition}

%%%%%%%%%%%%%%%%%%%%%%%%%%%%%%%%%%%%%%%%%%%%%%%%%%%%%%%%%%%%%%%%%%%%%%%%%%%%%%%
\subsection{Summary: Theory Space}
\label{app:theoryspace}
%%%%%%%%%%%%%%%%%%%%%%%%%%%%%%%%%%%%%%%%%%%%%%%%%%%%%%%%%%%%%%%%%%%%%%%%%%%%%%%

\begin{center}
\renewcommand{\arraystretch}{1.3}
\small
\begin{tabular}{lcccccc}
\hline
\textbf{Theory} & $C_{ij}$ & $\Psi$ dyn. & $U^{IDN}$ & $\lambda$ & Time & Level \\
\hline
Arrow--Debreu & $= 0$ & Exog. & No & $= 1$ & Static & Ind. \\
Fehr--Schmidt & $\neq 0$ & Slow & No & $= 1$ & Static & Ind. \\
Bolton--Ockenfels & $\neq 0$ & Slow & No & $= 1$ & Static & Ind. \\
Akerlof--Kranton & $\neq 0$ & Slow & \textbf{Yes} & $= 1$ & Dynamic & Ind. \\
Bénabou--Tirole & $\neq 0$ & Slow & \textbf{Yes} & $= 1$ & Dynamic & Ind. \\
Aghion--Howitt & $\neq 0$ & \textbf{Endog.} & No & $= 1$ & Dynamic & Macro \\
Acemoglu--Robinson & $\neq 0$ & \textbf{Endog.} & No & $= 1$ & Dynamic & Nation \\
Roberts & $> 0$ & Slow & No & $= 1$ & Static & Org. \\
Milgrom--Roberts & $\geq 0$ & Exog. & No & $= 1$ & Static & General \\
CAPM & $= \text{Cov}$ & Exog. & No & $= 1$ & Static & Market \\
Jensen--Meckling & $\lessgtr 0$ & Slow & No & $= 1$ & Static & Org. \\
Prospect Theory & $\neq 0$ & Endog. & No & $\mathbf{> 1}$ & Static & Ind. \\
\textbf{Becker--Goodman} & $\geq 0$ & Exog. & No & $= 1$ & \textbf{$\Psi_T$} & Ind./Market \\
\hline
\textbf{Full Framework} & Any & Endog. & Yes & $> 1$ & Dynamic & Multi \\
\hline
\end{tabular}
\end{center}

Each established theory is a \emph{restriction} of the full $(C, \Psi)$ framework.
The framework unifies them by relaxing different restrictions.

%==============================================================================
% PART 3: KEY RESULTS
%==============================================================================

% \section{Results} marker for template compliance

\subsection{Summary of Key Results}
\label{sec:A-results}

\begin{table}[htbp]
\centering
\caption{Summary of Formal Results in Appendix DER}
\label{tab:A-results}
\renewcommand{\arraystretch}{1.2}
\begin{tabular}{@{}llp{7cm}@{}}
\toprule
\textbf{Result} & \textbf{Type} & \textbf{Statement} \\
\midrule
A-R1 & Nesting & 12 major theories are limiting cases of $(C, \Psi)$ \\
A-R2 & Separability & Arrow-Debreu: $C_{ij} = 0 \Leftrightarrow$ no social preferences \\
A-R3 & Social Pref & Fehr-Schmidt/ERC: $C_{ij} \neq 0$ captures fairness \\
A-R4 & Identity & Akerlof-Kranton: $U^{IDN} \neq 0$ adds identity dimension \\
A-R5 & Dynamics & Aghion-Howitt/Acemoglu: Endogenous $\Psi$ for growth/institutions \\
\bottomrule
\end{tabular}
\end{table}

%==============================================================================
% PART 4: WORKED EXAMPLE
%==============================================================================

\subsection{Worked Example: From Arrow-Debreu to Fehr-Schmidt}
\label{sec:A-example}

\begin{tcolorbox}[colback=blue!5!white, colframe=blue!75!black,
    title=\textbf{Worked Example: Relaxing Separability}]

\textbf{Step 1: Start with Arrow-Debreu}

Standard model: $U_i(x_i)$ with $C_{ij} = 0$ for $i \neq j$.

Prediction: In ultimatum game, responders accept any positive offer.

\textbf{Step 2: Introduce Cross-Agent Dependence}

Allow $U_i = U_i(x_i, x_{-i})$ with:
\[
C_{ij} = \frac{\partial^2 U_i}{\partial x_i \partial x_j} \neq 0
\]

\textbf{Step 3: Specify Fehr-Schmidt Structure}

\[
U_i = x_i - \alpha_i \cdot \frac{1}{n-1} \sum_{j \neq i} \max(x_j - x_i, 0)
     - \beta_i \cdot \frac{1}{n-1} \sum_{j \neq i} \max(x_i - x_j, 0)
\]

This gives $C_{ij} = \frac{\alpha_i + \beta_i}{n-1} \cdot \delta(x_i - x_j)$.

\textbf{Step 4: Empirical Test}

With $\alpha = 2$, $\beta = 0.6$: Responders reject offers $< 30\%$.

Matches ultimatum game data. The restriction $C_{ij} = 0$ fails empirically.

\textbf{Conclusion:} Fehr-Schmidt $\supset$ Arrow-Debreu via relaxing separability.

\end{tcolorbox}

%==============================================================================
% IMPLICATIONS
%==============================================================================

\subsection{Implications for Economic Theory}
\label{sec:A-implications}

\begin{itemize}[nosep]
    \item \textbf{Theory choice:} Different contexts call for different restrictions
    \item \textbf{Empirical testing:} Test which restrictions are violated
    \item \textbf{Unification:} No need to choose between theories---they coexist
    \item \textbf{Extension:} New theories emerge by relaxing new restrictions
\end{itemize}

%==============================================================================
% PART 5: BACK MATTER
%==============================================================================

%------------------------------------------------------------------------------
% SUMMARY
%------------------------------------------------------------------------------

\subsection{Summary}
\label{sec:A-summary}

\begin{tcolorbox}[colback=blue!5!white, colframe=blue!75!black,
    title=\textbf{Summary: Limiting Cases}]

\textbf{What We Showed:}
\begin{itemize}[nosep]
    \item 12 major economic theories are special cases of $(C, \Psi)$
    \item Key dimensions: $C_{ij}$, $\Psi$ dynamics, $U^{IDN}$, $\lambda$
    \item Each theory = specific restrictions on general framework
    \item EBF unifies rather than replaces existing theories
\end{itemize}

\textbf{Theoretical Contribution:}
Provides formal foundation for interdisciplinary integration---behavioral,
institutional, financial, and organizational economics share common structure.

\end{tcolorbox}

%------------------------------------------------------------------------------
% GLOSSARY OF SYMBOLS
%------------------------------------------------------------------------------

\subsection{Glossary of Symbols}
\label{sec:A-glossary}

\begin{center}
\begin{tabular}{@{}lll@{}}
\toprule
\textbf{Symbol} & \textbf{Meaning} & \textbf{Reference} \\
\midrule
$C_{ij}$ & Cross-partial derivative & Eq. \ref{eq:generalcomplementarity} \\
$\Psi$ & Context vector & Eq. \ref{eq:generalcontext} \\
$U^{IDN}$ & Identity utility component & Section \ref{app:akerlofkranton} \\
$\lambda$ & Loss aversion parameter & Section \ref{app:prospecttheory} \\
$K$ & Coherence index & Throughout \\
$Q$ & Quality/welfare measure & Throughout \\
$\alpha_i, \beta_i$ & Inequity aversion parameters & Eq. \ref{eq:fsutility} \\
$C^*$ & Reference complementarity structure & Throughout \\
\bottomrule
\end{tabular}
\end{center}

\textbf{See also:} Appendix GLS (Central Glossary and Notation) for complete symbol definitions.

%------------------------------------------------------------------------------
% CRITICAL FOUNDATIONS
%------------------------------------------------------------------------------

\subsection{Critical Foundations}
\label{sec:A-foundations}

\begin{tcolorbox}[colback=red!5!white, colframe=red!75!black,
    title=\textbf{Critical Foundations: Objections and Responses}]

\textbf{Objection 1: The unification is trivial---any framework can nest others.}

\textit{``With enough parameters, you can fit anything.''}

\textbf{Response:} The nesting is \textit{structural}, not parametric. Each theory
corresponds to setting $C_{ij}$ to specific functional forms, not just fitting
free parameters. The restrictions are testable and economically meaningful.

\medskip

\textbf{Objection 2: Some theories are incompatible.}

\textit{``Arrow-Debreu assumes rationality; behavioral economics rejects it.''}

\textbf{Response:} The framework accommodates both: Arrow-Debreu as $C_{ij} = 0$
is a \textit{special case} that holds approximately in some contexts. Behavioral
deviations arise when $C_{ij} \neq 0$. The question becomes empirical: which
restrictions hold in which contexts?

\medskip

\textbf{Objection 3: The framework adds nothing new.}

\textit{``This is just relabeling existing concepts.''}

\textbf{Response:} The value is in the \textit{unified notation} and \textit{comparative
statics across theories}. By expressing all theories in $(C, \Psi)$ language,
we can: (1) identify which restrictions drive results; (2) systematically explore
intermediate cases; (3) generate novel predictions by relaxing multiple restrictions.

\end{tcolorbox}

%------------------------------------------------------------------------------
% OPEN ISSUES
%------------------------------------------------------------------------------

\subsection{Open Issues and Research Agenda}
\label{sec:A-openissues}

\paragraph{Issue 1: Additional Theories.}
This appendix covers 12 theories. Others (matching theory, mechanism design,
network economics) could be added to complete the mapping.

\paragraph{Issue 2: Empirical Discrimination.}
Which restrictions are most frequently violated empirically? A systematic
meta-analysis could identify the ``typical'' configuration across domains.

\paragraph{Issue 3: Computational Implementation.}
Translating the theoretical nesting into computational models would enable
quantitative comparison of predictions across frameworks.

\paragraph{Issue 4: Cross-Domain Transfer.}
If a restriction fails in one domain, does it fail similarly in others?
This could identify universal vs. domain-specific aspects of behavior.

%------------------------------------------------------------------------------
% REFERENCES
%------------------------------------------------------------------------------

\subsection{References for Appendix DER}
\label{sec:A-references}

\begin{itemize}[nosep]
    \item \citet{arrowdeb}: General equilibrium foundation
    \item \citet{fehrschmidt}: Inequity aversion model
    \item \citet{boltonockenfels}: ERC social preferences
    \item \citet{akerlofkranton}: Identity economics
    \item \citet{benaboutirole}: Moral identity model
    \item \citet{aghionhowitt}: Endogenous growth
    \item \citet{acemoglu}: Institutional economics
    \item \citet{roberts2004}: Organizational complementarities
    \item \citet{milgromroberts1990}: Strategic complementarity
    \item \citet{kahnemantversky1979}: Prospect theory
\end{itemize}

\textbf{Full citations:} See Master Bibliography (\texttt{bibliography/bcm\_master.bib})

%------------------------------------------------------------------------------
% CROSS-REFERENCE FOOTER
%------------------------------------------------------------------------------

\begin{tcolorbox}[colback=gray!10!white, colframe=gray!75!black,
    title=\textbf{Cross-Reference Footer}]
\label{sec:A-footer}

\textbf{This Appendix (A) connects to:}
\begin{itemize}[nosep]
    \item \textbf{B} (Complementarity): Detailed $C_{ij}$ structure
    \item \textbf{D} (Proofs): Formal derivations of $C^*$
    \item \textbf{I-O} (Literature): Empirical support for each theory
    \item \textbf{AA-AG} (Domains): Applied limiting cases
\end{itemize}

\textbf{Reading Path:}
$\rightarrow$ Chapter 5 (intuition) $\rightarrow$ This Appendix (formal)
$\rightarrow$ Appendix HOW (complementarity details) $\rightarrow$ Appendix PRF (proofs)

\end{tcolorbox}
